\documentclass[runningheads,envcountsect,envcountsame]{llncs}

\usepackage{amsmath,amssymb}
\usepackage[expansion=false]{microtype}
\usepackage{booktabs}
\usepackage{xcolor}
\usepackage{hyperref}
\renewcommand\UrlFont{\color{blue}\rmfamily}


\newcommand{\Mov}{\mathrm{Mov}}
\newcommand{\Det}{\mathrm{Det}}
\newcommand{\DetW}{\mathrm{Det}_{\omega}}
\newcommand{\Foot}{\mathrm{Foot}}
\newcommand{\score}{\mathrm{score}}
\newcommand{\TR}{\mathcal{T}(R)}

\begin{document}

\title{What a Mutation Score Certifies}
\subtitle{Adequacy-Completeness for Output Mutation}
\titlerunning{What a Mutation Score Certifies}

\author{Author Name}
\authorrunning{A. Author}
\institute{Affiliation, City, Country \\ \email{author@example.org}}

\maketitle

\begin{abstract}
Mutation testing scores a test suite by the fraction of seeded faults it
distinguishes, and a score of $1$ is read as evidence of adequacy. \emph{Adequate
for what} has never been answered: the field's formal results characterise
operator sets relative to a fixed program text or a fixed test suite, and its
bridge from seeded faults to real ones rests on the coupling-effect hypothesis.
We give a completeness notion that is neither program- nor suite-relative, for
the class of \emph{output mutation} operators (perturb the returned value;
extreme or Descartes-style mutation is the constant special case). For adequacy
such an operator is exactly its \emph{footprint}, the set of values it changes.
We call a family $\Pi$ adequacy-complete for a target family $\Gamma$ when, for
every program and every suite, killing all non-equivalent $\Pi$-mutants forces
killing all non-equivalent $\Gamma$-mutants, and prove an exact
characterisation: completeness holds iff every target footprint is the union of
the $\Pi$-footprints it contains. Consequences: the minimal absolutely complete
family on $n$ values consists of the $n$ singleton \emph{value guards}, and none
is finite on an infinite return type; a full absolute score certifies exactly
output coverage and nothing more; first-order completeness does not entail
higher-order completeness; completeness is decidable exactly when footprint
containment is; and the minimum certifying suite has size $|f(D)|$. The
development relativises to a weak oracle, yielding a precise characterisation of
pseudo-tested methods that an empirical study of five Python libraries bears out:
$43\%$ of covered methods are pseudo-tested under a value oracle against only
$6\%$ under a crash-inclusive one. The core is mechanised in Lean~4 / Mathlib.

\keywords{Mutation testing \and Test adequacy \and Operator completeness \and
Extreme mutation \and Interactive theorem proving}
\end{abstract}

\section{Introduction}
\label{sec:intro}

Mutation testing~\cite{demillo1978,hamlet1977} seeds a program with small faults
---\emph{mutants}---and measures a test suite by the fraction it \emph{kills},
i.e.\ distinguishes from the original. A mutation score of $1$ is the field's
adequacy certificate, and it is used as one: tools report it, teams gate on it,
and a growing body of work uses it to judge machine-generated code.

The certificate's strength rests entirely on the operator family. A score of $1$
says that no fault \emph{the operators can express} survived, and says nothing
about faults they cannot express. Two questions follow, and neither has a clean
answer in the literature.

\paragraph{What does a score of 1 certify?}
Not correctness. The bridge from ``no seeded fault survives'' to ``no real fault
survives'' is the \emph{coupling-effect hypothesis}~\cite{demillo1978} and the
\emph{competent-programmer hypothesis}, both stated as motivating assumptions and
supported empirically~\cite{offutt1992,jia2011,just2014} rather than proven.

\paragraph{When is an operator family complete?}
The strongest available results are relative by construction. \emph{Sufficient}
operator sets~\cite{offutt1996} are determined experimentally against a chosen
operator pool. \emph{Minimal} mutant sets and \emph{dominator}
mutants~\cite{ammann2014} are minimal relative to a fixed test set and an
already-generated pool. \emph{True} (semantic) mutant subsumption---the order
against which a completeness theorem would rank---is undecidable~\cite{kurtz2015},
and the foundational treatment of mutation adequacy~\cite{budd1982} likewise
establishes undecidability of the associated decision problems rather than a
characterisation.

The relativity is not an oversight. A program-text mutation \emph{schema} carries
no operator-only invariant that fixes its detection uniformly across programs:
the mutant it produces depends on the \emph{source text} of the program, not on
the function the program computes, so ``does schema $X$ subsume schema $Y$'' has
no program-independent answer (Section~\ref{sec:progindep}). Any completeness
notion for program-text operators must therefore quantify over programs
\emph{and} their texts, which is what forces the prior results to be
relative-by-construction.

\subsection{Approach}

We restrict attention to a sub-class that does admit an operator-only invariant:
\textbf{output mutation}, operators that perturb the value a function returns,
independent of how the function computes it. Extreme mutation---replacing a
method body with \texttt{return c}, as implemented in
Descartes~\cite{niedermayr2016,veraperez2018}---is the constant special case.

For this class we define completeness by \emph{what a kill certifies}, quantified
over all programs and all suites:

\begin{quote}
A family $\Pi$ is \textbf{adequacy-complete for a target family $\Gamma$} iff,
for every program and every finite suite, a mutation score of $1$ against $\Pi$
forces a mutation score of $1$ against $\Gamma$.
\end{quote}

Both relativities of the prior results are quantified away by the definition. The
resulting theory is governed by a single combinatorial object---the
\emph{footprint} of an operator, the set of values it changes---and reduces to
set containment.

\subsection{Contributions}

\begin{enumerate}
\item \textbf{The footprint reduction} (Section~\ref{sec:model}): for adequacy
purposes an output operator \emph{is} its footprint; equivalence and killing
depend only on which values the operator moves, never on where it sends them.

\item \textbf{A characterisation theorem} (Section~\ref{sec:char}): $\Pi$ is
adequacy-complete for $\Gamma$ iff every target footprint is the union of the
$\Pi$-footprints contained in it. We give the theorem for finite $\Pi$
(mechanised) and in a general form covering infinite families.

\item \textbf{The value-guard basis and the finite/infinite dichotomy}
(Section~\ref{sec:absolute}): on $n$ values the minimal absolutely complete family
has size exactly $n$, consisting of the singleton \emph{value guards}; on an
infinite return type no finite output-operator family is absolutely complete.
Descartes-style constants are complete iff $n = 2$.

\item \textbf{The ceiling} (Section~\ref{sec:absolute}): a full absolute score
holds exactly when the suite exercises every reachable output. This is the whole
operational content of output mutation, and it delimits the technique as much as
it licenses it.

\item \textbf{Relative completeness and its decidability}
(Section~\ref{sec:relative}): completeness reduces to footprint containment,
giving a P / decidable / undecidable spectrum that runs through set containment
rather than through algebraic word problems.

\item \textbf{Higher-order mutants} (Section~\ref{sec:ho}): first-order
completeness does not entail higher-order completeness; the only positive
statement is degenerate.

\item \textbf{Program independence} (Section~\ref{sec:progindep}):
output-mutant detection factors through the footprint, which is exactly why a
program-independent completeness notion is available here and not for
program-text schemas.

\item \textbf{The minimum certifying suite} (Section~\ref{sec:suite}): the
cheapest certifying suite has size $|f(D)|$, the teaching dimension of the
program against the operator family.

\item \textbf{Relativisation to a weak oracle} (Section~\ref{sec:oracle}): the
entire development goes through with the footprint replaced by an
oracle-relative detection set, yielding a precise characterisation of
pseudo-tested methods.
\end{enumerate}

The results of Sections~\ref{sec:model}--\ref{sec:suite} and the reduction from
programs and finite suites down to footprints are mechanised in Lean~4 /
Mathlib. Section~\ref{sec:mech} states exactly which results are machine-checked
and which are paper proofs.

\subsection{What this paper does not claim}
\label{sec:notclaim}

The account is bounded, and the bounds are load-bearing rather than decorative.
It concerns output mutation of pure total functions; it does not extend to
general program-text mutation, and we make no claim about the coupling hypothesis
for program-text faults. Sections~\ref{sec:scope} and~\ref{sec:threats} state the
model's assumptions and their consequences in full. The central quantitative
result---that a full absolute score is equivalent to output coverage---is
deliberately modest, and it implies that a tool gains no information from running
an absolutely complete output-operator family rather than measuring output
coverage directly. The contribution is a characterisation of what such a score
means, not a recommendation to compute it.

\section{Background and Related Work}
\label{sec:related}

\paragraph{Mutation testing.}
The technique originates with~\cite{demillo1978} and~\cite{hamlet1977};
\cite{jia2011} and~\cite{papadakis2019} survey its development. Its justification
rests on two hypotheses: the \emph{competent-programmer hypothesis}, that real
faults are close to correct programs, and the \emph{coupling effect}, that suites
detecting simple faults tend to detect complex ones. The coupling effect has
substantial empirical support~\cite{offutt1992,just2014} and no general proof.

\paragraph{Formal treatments.}
Budd and Angluin~\cite{budd1982} give the foundational formal analysis of
mutation adequacy, relating notions of correctness to testing and establishing
undecidability results for the associated decision problems. Their results are
undecidability theorems about a general setting, not a characterisation of which
operator families make a passing score adequate; the present paper is
complementary, giving such a characterisation for a restricted class where it is
available.

\paragraph{Operator selection.}
Offutt et al.~\cite{offutt1996} identify \emph{sufficient} operator sets
experimentally, by measuring which subsets of a fixed pool retain the pool's
discriminating power on a corpus. Ammann et al.~\cite{ammann2014} formalise
\emph{minimal} mutant sets and \emph{dominator} mutants: given a fixed program
and a fixed test pool, the dominator mutants are those whose kill implies the
kill of others. Kurtz et al.~\cite{kurtz2015} study mutant subsumption and show
that the true (semantic) subsumption relation is undecidable, restricting
practice to computable approximations. All three are relative to a fixed program,
a fixed pool, or both---which, as Section~\ref{sec:progindep} argues, is forced by
the absence of an operator-only invariant for program-text schemas.

\paragraph{Extreme mutation.}
Niedermayr et al.~\cite{niedermayr2016} introduce extreme mutation and the notion
of a \emph{pseudo-tested} method: one that is covered by the suite yet whose body
can be replaced wholesale without any test failing. Vera-P\'erez et
al.~\cite{veraperez2018} study the phenomenon at scale and provide the Descartes
tool. In the terms of this paper, extreme mutation is the constant
output-operator family; Section~\ref{sec:absolute} recovers it as exactly optimal
on two-valued return types and provably deficient beyond, and
Section~\ref{sec:oracle} gives a characterisation of pseudo-testedness in the
model's own terms.

\paragraph{Teaching dimension.}
The minimum number of labelled examples that identifies a concept within a class
is the \emph{teaching dimension}~\cite{goldman1995}. Section~\ref{sec:suite}
identifies the minimum certifying suite with this quantity for the
output-mutation setting.

\paragraph{Mechanisation.}
The development is checked in Lean~4~\cite{lean4} against Mathlib~\cite{mathlib}.

\section{The Model}
\label{sec:model}

\subsection{Programs, Suites, Operators}

\begin{definition}[program, suite, output operator]
\label{def:program}
A program computes a total function---its \emph{denotation}---$f : D \to R$, from
an input domain $D$ to a return type $R$ with $|R| \ge 2$. A \emph{test suite} is
a finite $T \subseteq D$. An \emph{output operator} is a total map $p : R \to R$;
it produces the \emph{mutant} $p \circ f$ (run the program, then apply $p$ to the
result). An \emph{output-operator family} $\Pi$ is a set of output operators.
\end{definition}

Post-composition is what an output or extreme mutation performs: it perturbs the
returned value uniformly, without reference to how the value was computed.

\begin{definition}[footprint]
\label{def:footprint}
The \emph{footprint} of an output operator $p$ is
\[
  \Mov(p) \;=\; \{\, r \in R : p(r) \neq r \,\} \;\subseteq\; R,
\]
the set of values $p$ actually changes. Write $I = f(D)$ for the program's
\textbf{reachable} outputs and $O = f(T)$ for the suite's \textbf{observed}
outputs; note $O \subseteq I$.
\end{definition}

\begin{definition}[equivalence, killing, score]
\label{def:score}
For a program $f$ and suite $T$:
\begin{itemize}
\item $p \circ f$ is \textbf{equivalent} to $f$ iff $p(f(x)) = f(x)$ for all
$x \in D$;
\item $T$ \textbf{kills} $p \circ f$ iff $p(f(x)) \neq f(x)$ for some $x \in T$;
\item the \textbf{score} $\score_\Pi(f,T) = 1$ iff $T$ kills every non-equivalent
mutant $p \circ f$, $p \in \Pi$.
\end{itemize}
\end{definition}

Equivalent mutants are excluded from the score, as in standard practice: they
cannot be killed by any suite, so counting them would make a score of $1$
unachievable.

\subsection{The Footprint Reduction}

\begin{proposition}[an operator is its footprint]
\label{prop:reduction}
For any program $f$, suite $T$, and operator $p$:
\[
  p \circ f \text{ is equivalent} \iff \Mov(p) \cap I = \varnothing,
  \qquad
  T \text{ kills } p \circ f \iff \Mov(p) \cap O \neq \varnothing.
\]
Consequently, for fixed $(f,T)$, whether $p$ is a non-equivalent survivor depends
\textbf{only} on $\Mov(p)$; two operators with equal footprints are
interchangeable; and $\score_\Pi(f,T)$ depends only on the set of footprints
$\Foot(\Pi) = \{\Mov(p) : p \in \Pi\}$.
\end{proposition}

\begin{proof}
$p \circ f$ is equivalent iff $p$ fixes $f(x)$ for every $x \in D$, i.e.\ $p$
fixes every element of $I = f(D)$, i.e.\ no element of $I$ is moved. $T$ kills
$p \circ f$ iff some $x \in T$ has $p(f(x)) \neq f(x)$, i.e.\ $f(x) \in \Mov(p)$
for some $x \in T$, i.e.\ $\Mov(p)$ meets $f(T) = O$. Both conditions read only
$\Mov(p)$ against the fixed sets $I$ and $O$. \qed
\end{proof}

Proposition~\ref{prop:reduction} is the engine of the paper. Where an operator
\emph{sends} the values it moves is invisible to adequacy; only \emph{which}
values it moves matters.

\begin{proposition}[every subset is a footprint]
\label{prop:surj}
For $|R| \ge 2$, every $S \subseteq R$ is the footprint of some output operator.
\end{proposition}

\begin{proof}
Fix distinct $a \neq b$. Let $q(a) = b$ and $q(x) = a$ for $x \neq a$; then $q$ is
fixed-point-free. Define $p(x) = q(x)$ for $x \in S$ and $p(x) = x$ otherwise.
Then $\Mov(p) = S$. \qed
\end{proof}

The construction avoids cycles: a single $|R|$-cycle would not serve for
uncountable $R$, whose orbits under iteration are countable. So footprints range
over the entire powerset $2^R$, and by Proposition~\ref{prop:reduction} the
adequacy content of a family $\Pi$ is exactly $\Foot(\Pi)$.

\subsection{Notation}

Throughout, $\Pi$ denotes the operator family under test and $\Gamma$ the target
family whose faults we wish to certify against; $\TR$ denotes the family of
\emph{all} output operators on $R$, and $n = |R|$ when $R$ is finite.

\subsection{Scope and Assumptions}
\label{sec:scope}

The model makes three assumptions. One is load-bearing and two are simplifying;
we state them here, where the model is defined, and return to their consequences
in Section~\ref{sec:threats}.

\paragraph{Load-bearing: the exact oracle.}
Definition~\ref{def:score} kills a mutant exactly when some test yields a
different return value, which presumes the suite observes the full output. Under
a weaker oracle a covered mutant can survive because no assertion inspects the
changed value---the \emph{pseudo-tested method}
phenomenon~\cite{niedermayr2016,veraperez2018} that motivates extreme mutation in
practice. Section~\ref{sec:oracle} relativises the entire development to an
arbitrary oracle and recovers the exact-oracle case as the special case where the
oracle is injective; readers who care primarily about weak oracles may read
Section~\ref{sec:oracle} as the general theory and
Sections~\ref{sec:char}--\ref{sec:suite} as its cleanest instance.

\paragraph{Simplifying: pure total functions.}
A denotation is a total map $f : D \to R$. Void methods, exceptions,
non-termination, and side effects are outside the model. Extending to partial or
effectful denotations requires enlarging $R$ with the relevant observations (an
exception tag, an effect trace), after which the footprint machinery applies
unchanged to the enlarged type; we do not develop this.

\paragraph{Simplifying: uniform post-composition.}
An output operator perturbs every returned value identically, $f \mapsto p \circ
f$. A statement-level mutation of a single \texttt{return} site is not of this
form.

\section{Adequacy-Completeness and the Characterisation Theorem}
\label{sec:char}

\begin{definition}[adequacy-completeness]
\label{def:complete}
Let $\Pi$ and $\Gamma$ be output-operator families. $\Pi$ is
\textbf{adequacy-complete for $\Gamma$} iff for \emph{every} program
$f : D \to R$ and \emph{every} finite suite $T \subseteq D$,
\[
  \score_\Pi(f,T) = 1 \;\Longrightarrow\; \score_\Gamma(f,T) = 1.
\]
$\Pi$ is \textbf{absolutely complete} iff it is complete for $\Gamma = \TR$.
\end{definition}

By Proposition~\ref{prop:reduction} this depends only on $\Foot(\Pi)$ and
$\Foot(\Gamma)$. The quantifier over all $(f,T)$ removes both the program-text
relativity and the test-suite relativity of prior completeness notions: the
definition asks what a kill certifies \emph{as such}, not what it certifies on
some corpus.

\begin{theorem}[footprint characterisation, finite $\Pi$]
\label{thm:char}
Let $\Pi$ be finite. Then $\Pi$ is adequacy-complete for $\Gamma$ iff
\[
  \forall g \in \Gamma,\ \forall b \in \Mov(g),\ \exists p \in \Pi :\quad
  b \in \Mov(p) \subseteq \Mov(g).
\]
Equivalently: \textbf{every target footprint is the union of the
$\Pi$-footprints it contains},
$\Mov(g) = \bigcup \{\Mov(p) : p \in \Pi,\ \Mov(p) \subseteq \Mov(g)\}$.
\end{theorem}

\begin{proof}
($\Leftarrow$) Assume the condition. Fix $(f,T)$ with $\score_\Pi(f,T) = 1$ and a
non-equivalent $g \in \Gamma$; we show $T$ kills $g \circ f$. Non-equivalence
gives (Proposition~\ref{prop:reduction}) $\Mov(g) \cap I \neq \varnothing$; pick
$b \in \Mov(g) \cap I$. By hypothesis there is $p \in \Pi$ with
$b \in \Mov(p) \subseteq \Mov(g)$. Then $b \in \Mov(p) \cap I$, so $p \circ f$ is
non-equivalent, so---the $\Pi$-score being $1$---it is killed:
$\Mov(p) \cap O \neq \varnothing$. Since $\Mov(p) \subseteq \Mov(g)$, also
$\Mov(g) \cap O \neq \varnothing$, i.e.\ $T$ kills $g \circ f$.

($\Rightarrow$) By contraposition. Assume the condition fails, witnessed by
$g \in \Gamma$ and $b \in \Mov(g)$ such that every $p \in \Pi$ with
$b \in \Mov(p)$ has an \emph{escape point} $a_p \in \Mov(p) \setminus \Mov(g)$.
Let $A = \{a_p : p \in \Pi,\ b \in \Mov(p)\}$. Since $\Pi$ is finite, $A$ is
finite, and $b \notin A$ (each $a_p$ avoids $\Mov(g)$, while $b \in \Mov(g)$).

Build the separating instance: let $I^\star = \{b\} \cup A$, take $D = I^\star$
and $f = \mathrm{id}$, so $I = I^\star$; take the suite $T = A$, so $O = A$. Then:

\emph{Every non-equivalent $\Pi$-mutant is killed.} Let $p \in \Pi$ be
non-equivalent, i.e.\ $\Mov(p) \cap I \neq \varnothing$. If $b \in \Mov(p)$, then
$a_p \in \Mov(p) \cap A = \Mov(p) \cap O$, so $p$ is killed. If
$b \notin \Mov(p)$, then
$\Mov(p) \cap I = \Mov(p) \cap (\{b\} \cup A) = \Mov(p) \cap A \neq \varnothing$,
so $p$ is killed. Hence $\score_\Pi(f,T) = 1$.

\emph{$g \circ f$ survives.} $b \in \Mov(g) \cap I$, so $g \circ f$ is
non-equivalent; but $\Mov(g) \cap O = \Mov(g) \cap A = \varnothing$, since every
$a_p$ avoids $\Mov(g)$. So $T$ does not kill $g \circ f$, and
$\score_\Gamma(f,T) < 1$. \qed
\end{proof}

Finiteness of $\Pi$ is used only to keep $A$ finite, so that the separating
program has a finite reachable set. The general case weakens the containment to
an avoidance condition.

\begin{theorem}[footprint characterisation, general $\Pi$]
\label{thm:chargen}
For arbitrary $\Pi$ and $\Gamma$, $\Pi$ is adequacy-complete for $\Gamma$ iff
\[
  \forall g \in \Gamma,\ \forall b \in \Mov(g),\
  \forall\, F \subseteq R \setminus \Mov(g) \text{ finite},\
  \exists p \in \Pi :\ b \in \Mov(p) \ \wedge\ \Mov(p) \cap F = \varnothing.
\]
\end{theorem}

\begin{proof}
($\Leftarrow$) Assume the condition and fix $(f,T)$ with $\score_\Pi(f,T) = 1$
and non-equivalent $g \in \Gamma$. Pick $b \in \Mov(g) \cap I$. Let
$F = O \setminus \Mov(g)$, which is finite (as $O$ is) and disjoint from
$\Mov(g)$. Obtain $p \in \Pi$ with $b \in \Mov(p)$ and
$\Mov(p) \cap F = \varnothing$. Since $b \in \Mov(p) \cap I$, $p \circ f$ is
non-equivalent, hence killed: there is $c \in O \cap \Mov(p)$. Since $\Mov(p)$
avoids $F = O \setminus \Mov(g)$ and $c \in O$, we get $c \in \Mov(g)$. So
$\Mov(g) \cap O \neq \varnothing$ and $T$ kills $g \circ f$.

($\Rightarrow$) By contraposition. Suppose the condition fails: there are
$g \in \Gamma$, $b \in \Mov(g)$, and a finite $F \subseteq R \setminus \Mov(g)$
such that every $p \in \Pi$ with $b \in \Mov(p)$ meets $F$. Take
$I = \{b\} \cup F$, $D = I$, $f = \mathrm{id}$, $T = F$, so $O = F$. If
$p \in \Pi$ is non-equivalent and $b \in \Mov(p)$, then
$\Mov(p) \cap F \neq \varnothing$ by assumption, so $p$ is killed; if
$b \notin \Mov(p)$, then $\Mov(p) \cap I = \Mov(p) \cap F \neq \varnothing$, so
$p$ is killed. Hence $\score_\Pi(f,T) = 1$. But $b \in \Mov(g) \cap I$ while
$\Mov(g) \cap O = \Mov(g) \cap F = \varnothing$, so $g \circ f$ is a
non-equivalent survivor. \qed
\end{proof}

\begin{corollary}
\label{cor:genfin}
Theorem~\ref{thm:char} is the case of Theorem~\ref{thm:chargen} for finite $\Pi$.
\end{corollary}

\begin{proof}
If the condition of Theorem~\ref{thm:char} holds, then for any finite $F$
disjoint from $\Mov(g)$ the witness $p$ with $\Mov(p) \subseteq \Mov(g)$
satisfies $\Mov(p) \cap F = \varnothing$. Conversely, suppose the condition of
Theorem~\ref{thm:chargen} holds and $\Pi$ is finite. If no $p$ had
$b \in \Mov(p) \subseteq \Mov(g)$, then each $p$ with $b \in \Mov(p)$ would have
an escape point $a_p \notin \Mov(g)$; the set $A$ of these is finite and disjoint
from $\Mov(g)$, so applying Theorem~\ref{thm:chargen} with $F = A$ yields a $p$
with $b \in \Mov(p)$ and $\Mov(p) \cap A = \varnothing$, contradicting
$a_p \in \Mov(p) \cap A$. \qed
\end{proof}

Theorem~\ref{thm:chargen} matters because several natural families are
infinite---for instance the full value-guard family on an infinite return type
(Section~\ref{sec:absolute}), which is absolutely complete but not finite, and
therefore outside the reach of Theorem~\ref{thm:char}.

\begin{remark}[composition is the wrong closure]
\label{rem:composition}
Theorem~\ref{thm:char} dissolves the temptation to define completeness by
\emph{generation under composition}. Killing the operators of $\Pi$ does not kill
their composites. Take $R = \{0,1,2\}$, a program with sole reachable and
observed value $0$, and operators $a,b$ with $a(0) = 1$, $b(0) = 2$, $b(1) = 0$.
Both $a \circ f$ and $b \circ f$ are killed at $0$, yet $(b \circ a)(0) = b(1) =
0$, so $(b \circ a) \circ f$ is not killed---it is equivalent on this program.
Composition acts precisely on the part of an operator that adequacy cannot see:
$\Mov(b \circ a)$ can be strictly smaller than both $\Mov(a)$ and $\Mov(b)$. So
``$\Pi$ generates the full transformation monoid'' says nothing about what a
score of $1$ certifies. The governing order is footprint containment, not
compositional generation.
\end{remark}

\section{Absolute Completeness: Value Guards, the Ceiling, and the Dichotomy}
\label{sec:absolute}

\begin{corollary}[the value-guard basis]
\label{cor:guards}
Let $|R| = n$. A family $\Pi$ is absolutely complete iff it contains, for every
$r \in R$, an operator whose footprint is the singleton $\{r\}$---a \emph{value
guard} $\gamma_r$, readable as ``if the result is $r$, return something else.''
The minimal absolutely complete family is $\{\gamma_r : r \in R\}$, of size
exactly $n$.
\end{corollary}

\begin{proof}
Absolute completeness is completeness for $\Gamma = \TR$, whose footprints are
all of $2^R$ (Proposition~\ref{prop:surj}). Apply Theorem~\ref{thm:char} to the
singleton target footprints $\{b\}$: completeness forces, for each $b$, some
$p \in \Pi$ with $b \in \Mov(p) \subseteq \{b\}$, i.e.\ $\Mov(p) = \{b\}$.
Conversely, given a guard for every value, any target footprint $S$ and any
$b \in S$ admit $\gamma_b$ with $b \in \{b\} \subseteq S$, so
Theorem~\ref{thm:char}'s condition holds. Minimality: each of the $n$ singleton
targets forces a distinct guard. \qed
\end{proof}

\begin{corollary}[the finite/infinite dichotomy]
\label{cor:infinite}
A \emph{finite} operator family can be absolutely complete only when $R$ is
finite. Equivalently: for an infinite return type---\texttt{int},
\texttt{String}, arbitrary objects---\textbf{no finite output-operator family is
absolutely complete.}
\end{corollary}

\begin{proof}
By Corollary~\ref{cor:guards} absolute completeness requires a distinct singleton
footprint $\{r\}$ for every $r \in R$; the map $r \mapsto \{r\}$ is injective, so
a complete family has at least $|R|$ members. \qed
\end{proof}

Corollary~\ref{cor:infinite} bounds the value-guard story to finite return
types---Booleans, enums, small tagged unions. On the return types of most real
programs the absolute notion is unavailable, and the meaningful question is
\emph{relative} completeness against a chosen target class
(Section~\ref{sec:relative}). That is where a certificate a tool can actually
offer lives.

\begin{corollary}[extreme mutation is optimal on Booleans, deficient beyond]
\label{cor:constants}
A Descartes-style constant operator \texttt{return c} has footprint
$R \setminus \{c\}$. The constant family $\{\texttt{return } c : c \in R\}$ is
absolutely complete iff $n = 2$.
\end{corollary}

\begin{proof}
By Corollary~\ref{cor:guards}, absolute completeness requires a singleton
footprint for each value; $|R \setminus \{c\}| = 1$ iff $n = 2$. For $n = 2$,
\texttt{return true} has footprint $\{\mathrm{false}\}$ and \texttt{return false}
has footprint $\{\mathrm{true}\}$---precisely the two value guards, so extreme
mutation \emph{is} the minimal complete family on Booleans. For $n \ge 3$ every
constant footprint has size at least $2$, so no constant is a value guard. \qed
\end{proof}

\begin{example}[a fault the constants miss]
\label{ex:miss}
Let $R = \{a,b,c\}$, a program reaching all three, and a suite observing only $a$
and $b$ (so $O = \{a,b\}$, $I = \{a,b,c\}$). All three constant mutants are
killed: \texttt{return a} moves $\{b,c\}$, observed at $b$; \texttt{return b}
moves $\{a,c\}$, observed at $a$; \texttt{return c} moves $\{a,b\}$, observed at
$a$. The constant score is $1$. But the fault ``returns $b$ where it should
return $c$''---an operator with footprint $\{c\}$---is non-equivalent ($c \in I$)
and unkilled ($\{c\} \cap O = \varnothing$). A perfect extreme-mutation score
certifies nothing about the $c$-outputs the suite never observed.
\end{example}

\begin{theorem}[the ceiling]
\label{thm:ceiling}
For $|R| \ge 2$ and any program $f$ and suite $T$:
\[
  \score_{\TR}(f,T) = 1 \iff I \subseteq O \iff O = I,
\]
i.e.\ an absolute mutation score of $1$ holds \textbf{exactly when the suite
exercises every reachable output}.
\end{theorem}

\begin{proof}
($\Leftarrow$) If $O = I$, every non-equivalent $g$ has
$\Mov(g) \cap I \neq \varnothing$, hence $\Mov(g) \cap O \neq \varnothing$, so
all die. ($\Rightarrow$) If $I \not\subseteq O$, pick $b \in I \setminus O$ and
the guard $\gamma_b$ with footprint $\{b\}$: it is non-equivalent ($b \in I$) and
unkilled ($\{b\} \cap O = \varnothing$), so the score is below $1$. \qed
\end{proof}

\begin{remark}[the ceiling, read operationally]
\label{rem:ceiling}
Theorem~\ref{thm:ceiling} gives the operational content of output mutation: a
full score certifies \emph{output coverage}---the dynamic property that every
value the program can return is observed by the suite---and the value guards are
the cheapest family making that certificate exact. It certifies nothing about
faults that separate two inputs mapped to the same output: if $f(x_1) = f(x_2)$
then $p(f(x_1)) = p(f(x_2))$ for every output operator $p$, so no output mutant
distinguishes $x_1$ from $x_2$. Input-separating faults are structurally outside
the reach of the technique.
\end{remark}

\section{Relative Completeness and Its Decidability}
\label{sec:relative}

Fix a finite $\Pi$ and a finitely presented target $\Gamma$. By
Theorem~\ref{thm:char}, deciding completeness is deciding a footprint-containment
condition, and its complexity is exactly that of the footprint representation.

\begin{proposition}[the exact reach of a partial guard family]
\label{prop:reach}
The value guards $\{\gamma_{r_1},\dots,\gamma_{r_k}\}$ are adequacy-complete for
exactly
$\Gamma_{\max} = \{\, g : \Mov(g) \subseteq \{r_1,\dots,r_k\} \,\}$.
\end{proposition}

\begin{proof}
By Theorem~\ref{thm:char} the guards are complete for $g$ iff every
$b \in \Mov(g)$ has a guard $\gamma_{r_i}$ with $b \in \{r_i\} \subseteq \Mov(g)$,
i.e.\ $b \in \{r_1,\dots,r_k\}$. This holds for all $b \in \Mov(g)$ iff
$\Mov(g) \subseteq \{r_1,\dots,r_k\}$. \qed
\end{proof}

Proposition~\ref{prop:reach} is the practically usable form of the theory. A
partial guard family yields an explicit, checkable certificate: \emph{every
guarded value the program can produce is exercised by the suite, hence any fault
confined to the guarded values is caught.} The guarded set is a parameter the
tool chooses and reports.

\begin{theorem}[decidability spectrum]
\label{thm:decide}
Deciding whether a finite $\Pi$ is adequacy-complete for a finite $\Gamma$
reduces to finitely many footprint-containment tests
$\Mov(p) \subseteq \Mov(g)$ and membership tests $b \in \Mov(p)$. Hence:
\begin{enumerate}
\item \textbf{Finite return type (value tables).} Footprints are explicit subsets
of an $n$-element set; each test is a linear scan, so completeness is decidable
in polynomial time.
\item \textbf{Semilinear footprints} ($R = \mathbb{Z}$, $\Mov$
Presburger-definable). Containment of Presburger-definable sets is decidable, so
completeness is decidable.
\item \textbf{Total computable operators.} Each footprint $\Mov(p)$ is a
\emph{decidable} set, but deciding containment---hence completeness---is
undecidable.
\end{enumerate}
\end{theorem}

\begin{proof}
The reduction is Theorem~\ref{thm:char}, a finite conjunction over $g \in \Gamma$
and $b \in \Mov(g)$, existentially quantifying $p \in \Pi$. (1) Tables give
constant-time membership and linear-time containment. (2) Presburger arithmetic
is decidable. (3) Reduce from non-halting: for a Turing machine $M$, the total
computable operator $g_M(k) = k+1$ if $M$ halts within $k$ steps and
$g_M(k) = k$ otherwise has decidable footprint
$\Mov(g_M) = \{k : M \text{ halts within } k \text{ steps}\}$, and
$\Mov(g_M) \subseteq \Mov(\mathrm{id}) = \varnothing$ iff $M$ never halts. \qed
\end{proof}

\begin{remark}[the spectrum runs through set containment]
\label{rem:spectrum}
The P / decidable / undecidable gradient here is that of subset containment for
increasingly expressive set representations. No monoid- or group-theoretic
machinery---word problems, submonoid membership, Cayley embeddings---appears,
because adequacy never composes operators (Remark~\ref{rem:composition}). It is
the containment order, not a generation order, that governs decidability. An
earlier formulation of operator completeness in terms of transformation-monoid
generation leads to exactly that algebraic machinery, and to a different and, for
adequacy purposes, incorrect answer; Remark~\ref{rem:composition} is the reason.
\end{remark}

\section{Higher-Order Mutants}
\label{sec:ho}

The coupling-effect hypothesis~\cite{demillo1978} is the empirical claim that
suites killing first-order mutants tend also to kill higher-order ones. For
higher-order \emph{output} mutants---compositions of output operators---no such
implication holds as a theorem, and the footprint lens shows exactly why.

\begin{proposition}[first-order completeness does not entail higher-order completeness]
\label{prop:coupling}
There exist a family $\Pi$, a program $f$, and a suite $T$ under which every
non-equivalent first-order $\Pi$-mutant is killed while a non-equivalent
higher-order mutant survives.
\end{proposition}

\begin{proof}
Let $R = \{0,1,2\}$ and $\Pi = \{a\}$ with
$a : 0 \mapsto 1,\ 1 \mapsto 0,\ 2 \mapsto 1$, so $\Mov(a) = R$. Take a program
$f$ with reachable set $I = \{0,2\}$ observed at $O = \{0\}$. The only
first-order mutant $a \circ f$ is non-equivalent and killed, since
$0 \in \Mov(a) \cap O$. The higher-order mutant $a \circ a$ satisfies
$a \circ a : 0 \mapsto 0,\ 1 \mapsto 1,\ 2 \mapsto 0$, so
$\Mov(a \circ a) = \{2\}$: it meets $I$ (at $2$, hence non-equivalent) and misses
$O$ (hence unkilled). So the score is $1$ against $\Pi$ and below $1$ against the
closure of $\Pi$ under composition. \qed
\end{proof}

By Proposition~\ref{prop:reduction} the reason is exact: killing sees only
$\Mov(\cdot) \cap O$, and $\Mov(a \circ a)$ can be strictly smaller than
$\Mov(a)$, removing precisely the value whose observation killed the first-order
mutant.

\begin{proposition}[the only positive statement is degenerate]
\label{prop:degenerate}
If $\Pi$ is absolutely complete it is adequacy-complete for \emph{every} target,
in particular for its closure under composition.
\end{proposition}

\begin{proof}
Corollary~\ref{cor:guards} gives $\Pi$ a value guard for every value;
Theorem~\ref{thm:char}'s condition then holds for any target. \qed
\end{proof}

Proposition~\ref{prop:degenerate} is a corollary of the value-guard basis, not a
coupling phenomenon: it holds because such a family already covers every
footprint, so composition is irrelevant.

\begin{remark}[what this does and does not say]
\label{rem:coupling}
By Theorem~\ref{thm:char}, higher-order-ness carries no special weight: a
higher-order mutant is simply an operator with a footprint, and the closure of
$\Pi$ under composition may contain footprints $\Pi$ does not cover. We therefore
claim only that coupling is not \emph{entailed} for output mutation---not that it
fails to occur. The coupling hypothesis is a statistical claim about typical
faults, which a single counterexample does not refute. Program-text coupling is a
separate question on which we make no claim.
\end{remark}

\section{Program Independence}
\label{sec:progindep}

\begin{proposition}[program-independent subsumption]
\label{prop:subsume}
For output operators $p$ and $g$: \emph{every test that kills $p \circ f$ also
kills $g \circ f$, for every program $f$}, iff $\Mov(p) \subseteq \Mov(g)$.
\end{proposition}

\begin{proof}
($\Leftarrow$) A test $x$ kills $p \circ f$ iff $f(x) \in \Mov(p) \subseteq
\Mov(g)$, so $x$ kills $g \circ f$. ($\Rightarrow$) If
$b \in \Mov(p) \setminus \Mov(g)$, take $f$ with $b$ reachable and a suite
observing $b$: it kills $p \circ f$ but not $g \circ f$. \qed
\end{proof}

So Theorem~\ref{thm:char} reads, in the language of mutant
subsumption~\cite{kurtz2015,ammann2014}: \textbf{$\Pi$ is adequacy-complete for
$\Gamma$ iff every target mutant is subsumed by some $\Pi$-mutant, uniformly over
all programs.} It is the program-independent form of dominator theory.

\begin{proposition}[the factoring that buys program independence]
\label{prop:factor}
For an output operator, the detection set factors through the footprint:
\[
  \Det(p \circ f) \;=\; \{x : p(f(x)) \neq f(x)\} \;=\; f^{-1}(\Mov(p)).
\]
So $p$ enters detection only through the program-independent object $\Mov(p)$:
the same $\Mov(p)$ governs detection across all programs $f$.
\end{proposition}

The contrast with program-text mutation is \emph{not} that its detection is
syntactic. The detection set $\Det(m) = \{x : m(x) \neq f(x)\}$ of any mutant is
semantic, depending only on the functions $m$ and $f$. The contrast is that a
program-text mutation \emph{schema} carries no operator-only invariant: the
mutant $\mu(\mathrm{source\ of\ } f)$ depends on the \emph{text} of $f$, so its
detection is a function of the (schema, source) pair rather than of $\mu$ alone,
and varies across source programs computing the same $f$. This is why the field's
completeness, sufficiency, and subsumption results are relative by construction
(Section~\ref{sec:related}), and why Definition~\ref{def:complete} is available
for output mutation and not for program-text mutation.

\section{The Minimum Certifying Suite}
\label{sec:suite}

\begin{theorem}[minimum certifying suite]
\label{thm:suite}
Fix an absolutely complete family and a program $f$ with reachable set $I =
f(D)$. A suite $T$ achieves $\score(f,T) = 1$ iff $O = f(T) = I$. If $I$ is
finite, the minimum certifying suite has size $\sigma(f) = |I| = |f(D)|$ (one
test per reachable value). If $I$ is infinite, no finite suite is certifying: a
finite $T$ has finite $O = f(T)$, so $O \neq I$.
\end{theorem}

\begin{proof}
By Theorem~\ref{thm:ceiling} the score is $1$ iff $O = I$. If $I$ is finite, a
suite with $O = I$ contains, for each $r \in I$, some $x$ with $f(x) = r$, so
$|T| \ge |I|$; choosing one such $x$ per value gives $|T| = |I|$. If $I$ is
infinite, $O = f(T)$ is finite for finite $T$. \qed
\end{proof}

\begin{remark}[sample complexity of certification]
\label{rem:teaching}
$\sigma(f)$ is the \emph{teaching dimension} of the program against the operator
family in the classical sense~\cite{goldman1995}: the minimum number of labelled
examples---here, tests together with their observed outputs---required to
distinguish $f$ from every non-equivalent mutant the family can express. For
output mutation it is exactly the number of reachable outputs. This is the
exact-identification companion to the completeness result:
Theorem~\ref{thm:char} says \emph{which} faults a family can certify against;
Theorem~\ref{thm:suite} says \emph{how many tests} the certification costs.
\end{remark}

\section{Beyond the Exact Oracle}
\label{sec:oracle}

The exact-oracle assumption (Section~\ref{sec:scope}) is the model's one
load-bearing simplification, and the one that most constrains practical
relevance, since the phenomenon motivating extreme mutation---pseudo-tested
methods~\cite{niedermayr2016,veraperez2018}---is a weak-oracle phenomenon by
definition. This section shows the development relativises, and that the
relativised model characterises pseudo-testedness.

\begin{definition}[oracle]
\label{def:oracle}
An \emph{oracle} is a map $\omega : R \to \Omega$ recording what the suite's
assertions actually observe about a returned value. The exact oracle is
$\omega = \mathrm{id}$. A suite that checks only a field, a hash, a type tag, or
nothing at all is a non-injective $\omega$; a constant $\omega$ models a suite
with no output assertions whatsoever.
\end{definition}

\begin{definition}[oracle-relative killing]
\label{def:okill}
$T$ \textbf{$\omega$-kills} $p \circ f$ iff $\omega(p(f(x))) \neq \omega(f(x))$
for some $x \in T$; $p \circ f$ is \textbf{$\omega$-equivalent} iff
$\omega(p(f(x))) = \omega(f(x))$ for all $x \in D$. The \textbf{$\omega$-score}
is $1$ iff every non-$\omega$-equivalent mutant is $\omega$-killed.
\end{definition}

Taking $\omega$-equivalence rather than true equivalence as the exclusion
criterion matches tool behaviour: a mutant no assertion can ever distinguish is
one the tool cannot kill and will report as equivalent.

\begin{definition}[oracle-relative detection set]
\label{def:detw}
$\DetW(p) = \{ r \in R : \omega(p(r)) \neq \omega(r) \}$.
\end{definition}

\begin{proposition}[the relativised reduction]
\label{prop:orelred}
$p \circ f$ is $\omega$-equivalent iff $\DetW(p) \cap I = \varnothing$; $T$
$\omega$-kills $p \circ f$ iff $\DetW(p) \cap O \neq \varnothing$. Moreover
$\DetW(p) \subseteq \Mov(p)$, with equality for injective $\omega$.
\end{proposition}

\begin{proof}
Identical to Proposition~\ref{prop:reduction}, reading $\omega \circ p$ against
$\omega$ in place of $p$ against the identity. Containment: if $p(r) = r$ then
$\omega(p(r)) = \omega(r)$. \qed
\end{proof}

\begin{corollary}[everything relativises]
\label{cor:relativise}
Since $\DetW(p) \subseteq R$ and Proposition~\ref{prop:orelred} has the form of
Proposition~\ref{prop:reduction}, Theorems~\ref{thm:char}, \ref{thm:chargen},
\ref{thm:ceiling}, \ref{thm:decide}, \ref{thm:suite} and
Propositions~\ref{prop:reach}, \ref{prop:coupling}, \ref{prop:degenerate},
\ref{prop:subsume} hold verbatim with $\Mov$ replaced throughout by $\DetW$.
\end{corollary}

The interesting content is where the two differ.

\begin{proposition}[the ceiling is oracle-independent for any non-trivial oracle]
\label{prop:oceiling}
If $\omega$ is non-constant, then for every $b \in R$ there is an operator with
$\DetW(p) = \{b\}$, and consequently the absolute $\omega$-score is $1$ iff
$I \subseteq O$.
\end{proposition}

\begin{proof}
Non-constant $\omega$ gives $r_1, r_2$ with $\omega(r_1) \neq \omega(r_2)$; for
any $b$, at least one of them, say $r'$, has $\omega(r') \neq \omega(b)$. Set
$p(b) = r'$ and $p = \mathrm{id}$ elsewhere; then $\DetW(p) = \{b\}$. Apply the
relativised Theorem~\ref{thm:ceiling}. \qed
\end{proof}

So weakening the oracle does not lower the ceiling: output coverage remains
exactly what a full absolute score certifies, provided the suite asserts
\emph{something}. What weakening the oracle changes is which \emph{specific}
families reach the ceiling---footprints shrink, so families lose discriminating
power.

\begin{proposition}[constants under a weak oracle]
\label{prop:oconst}
$\DetW(\texttt{return } c) = R \setminus \omega^{-1}(\omega(c))$. The constant
family is absolutely $\omega$-complete iff $|R| = 2$ and $\omega$ separates the
two values.
\end{proposition}

\begin{proof}
$\omega(c) \neq \omega(r)$ iff $r$ lies outside the $\omega$-class of $c$. By the
relativised Corollary~\ref{cor:guards}, completeness requires each
$\DetW(\texttt{return } c)$ to be a singleton for a suitable $c$ per value $b$,
i.e.\ $\omega^{-1}(\omega(c)) = R \setminus \{b\}$. For this to hold for every
$b$, the $\omega$-partition must be $\{\{b\}, R \setminus \{b\}\}$ for each $b$
simultaneously, which forces $|R| = 2$ with $\omega$ injective on $R$. \qed
\end{proof}

\begin{proposition}[pseudo-testedness, characterised]
\label{prop:pseudo}
Let $f$ be covered by $T$ (so $O \neq \varnothing$). The constant mutant
\texttt{return c} survives against $T$ under oracle $\omega$ while being
non-$\omega$-equivalent iff
\[
  O \subseteq \omega^{-1}(\omega(c)) \quad\text{and}\quad
  I \not\subseteq \omega^{-1}(\omega(c)),
\]
i.e.\ \textbf{every output the suite observes is oracle-indistinguishable from
$c$, while some reachable output is not.} Under the exact oracle this reduces to
$O = \{c\}$.
\end{proposition}

\begin{proof}
Survival is $\DetW(\texttt{return } c) \cap O = \varnothing$, i.e.\ $O$ avoids
$R \setminus \omega^{-1}(\omega(c))$, i.e.\ $O \subseteq \omega^{-1}(\omega(c))$.
Non-$\omega$-equivalence is $\DetW \cap I \neq \varnothing$, i.e.\
$I \not\subseteq \omega^{-1}(\omega(c))$. Under $\omega = \mathrm{id}$,
$\omega^{-1}(\omega(c)) = \{c\}$. \qed
\end{proof}

This is the model's account of the empirical phenomenon: a method is
pseudo-tested exactly when the suite's assertions collapse everything it observes
into a single oracle class, so replacing the body by a constant in that class
changes nothing the suite can see. It also explains why extreme mutation is a
productive \emph{detector} in practice even though the constant family is a weak
\emph{certifier} (Corollary~\ref{cor:constants}): the two roles are different,
and the gap $\Mov(p) \setminus \DetW(p)$ is precisely the oracle weakness the
technique surfaces.

\begin{remark}[status]
\label{rem:ostatus}
The oracle-relative detection set $\DetW$ and Propositions~\ref{prop:oceiling},
\ref{prop:oconst}, \ref{prop:pseudo} are mechanised directly (Section~\ref{sec:mech},
Table~\ref{tab:status}). What remains a paper proof is
Corollary~\ref{cor:relativise} itself: it is a substitution of one subset-valued
invariant for another, so mechanising the \emph{full} relativised theory is a matter of
abstracting the Lean development over the detection function rather than reproving it, and
we have not yet carried out that abstraction.
\end{remark}

\section{Empirical: Pseudo-Testedness Is Oracle-Relative}
\label{sec:empirical}

Proposition~\ref{prop:pseudo} predicts that whether a covered method is
\emph{pseudo-tested} is a property not of the method but of the oracle: the
constant mutant \texttt{return\,$c$} survives while remaining
non-$\omega$-equivalent exactly when every observed output is
$\omega$-indistinguishable from $c$. We test this directly, taking as two concrete
oracles the standard distinction between an \emph{assertion} kill---which pins the
returned value, the exact oracle $\omega = \mathrm{id}$---and a \emph{crash}
(exception or timeout) kill, which proves only that the mutant ran differently and
leaves the value unpinned: a weak, non-injective~$\omega$
(Definition~\ref{def:okill}).

\paragraph{Setup.}
We ran the constant perturbation---\texttt{return\,0} at each return site---against
the \emph{existing} test suites of five widely used Python libraries, over a seeded
random sample of public functions and methods. A method enters the denominator only
when a covering test actually executes the mutated return; an unexercised return is
untested, not pseudo-tested. Matching the whole-body \texttt{return\,$c$} of extreme
mutation~\cite{niedermayr2016,veraperez2018}, a multi-return method counts as
pseudo-tested under an oracle iff \emph{no} covered return is killed by that oracle.
We report, per project, the fraction pseudo-tested under the value oracle (no
assertion pins any return) and under the crash-inclusive oracle (no test---not even
a crash---kills any return). The measurement uses the same engine that realises the
paper's operators; a run is discarded unless its baseline passes and its trace is
complete.

\begin{table}[t]
\centering
\caption{Constant-mutant pseudo-testedness of covered functions/methods in five
Python libraries, under the value (assertion) oracle and the crash-inclusive
oracle. ($\dagger$)~\texttt{click} is shown but excluded from the pooled figure:
$48$ of its $60$ sampled methods time out under in-process profiling, leaving too
few measured to report.}
\label{tab:pseudo}
\begin{tabular}{@{}lrrr@{}}
\toprule
Project & Covered w/ const & Value oracle & Crash oracle \\
\midrule
boltons  & 27 & 63\% & 4\% \\
inflect  & 33 & 52\% & 9\% \\
toolz    & 46 & 33\% & 0\% \\
humanize & 16 & 19\% & 19\% \\
click$^\dagger$ & 2 & --- & --- \\
\midrule
\textbf{Pooled} & \textbf{122} & \textbf{43\%} & \textbf{6\%} \\
\bottomrule
\end{tabular}
\end{table}

\paragraph{Result.}
Across the $122$ covered functions and methods of the four measurable projects,
$43\%$ ($52/122$) are pseudo-tested under the value oracle---squarely within the
$1$--$46\%$ that Vera-P\'erez et al.\ measured across $21$ Java
projects~\cite{veraperez2018}---but only $6\%$ ($7/122$) under the crash-inclusive
oracle. The $37$-point gap is precisely the population of constant returns that a
\emph{traditional} mutation score counts as killed---a downstream crash caught
them---while the returned \emph{value} is unpinned. This is the phenomenon
Section~\ref{sec:oracle} formalises, measured directly: the same code is roughly
seven times more pseudo-tested once crash detection is discounted, so
pseudo-testedness is a property of the oracle and not of the suite alone,
confirming Proposition~\ref{prop:pseudo}. The effect is uneven and interpretable:
in \texttt{boltons} and \texttt{toolz} the surviving constants are almost all
crash-killed and value-unpinned ($63\%$ vs.\ $4\%$; $33\%$ vs.\ $0\%$), whereas
\texttt{humanize}'s string formatters, collapsed to a constant, are either
value-asserted or fully survive ($19\%$ under either oracle).

\paragraph{Threats to validity.}
The operator fixes $c = 0$; a family over several constants would only raise the
rate, so the reported figures are lower bounds on pseudo-testedness. The sample is
modest ($242$ probes) and, once class methods are included, skews toward accessors,
which are less value-asserted than free functions---\texttt{boltons} at $63\%$ sits
above the Java ceiling for this reason, so we read the oracle \emph{gap}, not the
absolute rate, as the finding. \texttt{click} reflects a harness rather than a
coverage limitation: its command/context objects do not profile in process. None of
these affects the qualitative result, which is the ratio between the two oracles on
identical code.

\section{Consequences for Mutation-Based Tools}
\label{sec:tools}

The results delimit as much as they prescribe. For a tool that certifies code
against an output or extreme mutation family:

\begin{itemize}
\item \textbf{Report a passing score as output coverage, and only that.} By
Theorem~\ref{thm:ceiling} an absolute score of $1$ says ``every value the program
can return was observed by the suite.'' Do not read it as evidence about faults
that separate inputs mapped to the same output (Remark~\ref{rem:ceiling}), nor
about faults a weaker oracle would miss (Section~\ref{sec:oracle}).

\item \textbf{Note that, for the absolutely complete family, output mutation adds
no information over coverage.} A value-guard score of $1$ is \emph{equivalent} to
output coverage, so a tool gains nothing by running the guard mutants rather than
measuring output coverage directly. The value of Theorem~\ref{thm:ceiling} is
that it says what the score means, not that it recommends computing it.

\item \textbf{Where the certificate is partial, state it explicitly}
(Proposition~\ref{prop:reach}). On a large or infinite return type absolute
completeness is unavailable (Corollary~\ref{cor:infinite}), so the useful output
is relative: choose a target class $\Gamma$, compute the footprint basis
Theorem~\ref{thm:char} requires for it, and report \emph{``complete against
$\Pi$, which is adequacy-complete for $\Gamma$; therefore this suite certifies
property $C$''} in place of an uninterpreted percentage. For a partial guard
family the certificate reads ``any fault confined to the guarded values is
caught,'' and the guarded set should be named.

\item \textbf{Prefer value guards to constants past Booleans}
(Corollary~\ref{cor:constants}), if running output mutants at all. Constants are
complete only for two-valued return types; beyond that they leave the surviving
faults of Example~\ref{ex:miss}.

\item \textbf{Do not rely on coupling for output mutation}
(Proposition~\ref{prop:coupling}). Killing first-order output mutants does not
entail killing their higher-order compositions; if those are in scope, a tool
must generate them.

\item \textbf{Report the certifying-suite size $\sigma(f) = |f(D)|$}
(Theorem~\ref{thm:suite}) as the cost and the content of the guarantee.
\end{itemize}

\section{Mechanisation}
\label{sec:mech}

The development is machine-checked in Lean~4 against Mathlib. The artifact
defines completeness \emph{directly over programs and finite
suites}---\texttt{ProgScore} and \texttt{ProgComplete} range over
$f : D \to R$ and \texttt{T : Finset D}---and then bridges to the footprint
level: \texttt{progScore\_iff\_scoreAt} (Proposition~\ref{prop:reduction})
reduces a program's score to the footprint-level \texttt{ScoreAt} at the
reachable set \texttt{range f} and observed set \texttt{f "" T};
\texttt{realizable} shows every reachable/observed pair $(I,O)$ with
$O \subseteq I$ is realised by an actual program;
\texttt{progComplete\_iff\_complete} glues these together; and
\texttt{progComplete\_characterization} restates the main theorem over programs.

Every mechanised theorem is \texttt{\#print axioms}-clean, depending on exactly
\texttt{[propext, Classical.choice, Quot.sound]} (\texttt{det\_factors} and
\texttt{detW\_const} on none), with no \texttt{sorryAx}.

\begin{table}[t]
\centering
\caption{Status of each result. ``Mechanised'' means a Lean theorem with no
\texttt{sorry} and a clean axiom audit.}
\label{tab:status}
\begin{tabular}{@{}llp{3.1cm}@{}}
\toprule
Result & Lean name & Status \\
\midrule
Prop.~\ref{prop:reduction} & \texttt{progScore\_iff\_scoreAt} & mechanised \\
Prop.~\ref{prop:surj} & \texttt{mov\_surjective}, \texttt{mov\_image\_univ} & mechanised \\
\textbf{Thm.~\ref{thm:char}} & \textbf{\texttt{footprint\_characterization}} & \textbf{mechanised} \\
Thm.~\ref{thm:char} (programs) & \texttt{progComplete\_characterization} & mechanised \\
Thm.~\ref{thm:chargen} & \texttt{footprint\_char\-acterization\_general} & mechanised \\
Cor.~\ref{cor:guards} & \texttt{absolute\_iff\_guards} & mechanised \\
Cor.~\ref{cor:infinite} & \texttt{progComplete\_univ\_infinite} & mechanised \\
Cor.~\ref{cor:constants} & \texttt{constants\_iff\_card\_two} & mechanised \\
Thm.~\ref{thm:ceiling} & \texttt{ceiling} & mechanised \\
Prop.~\ref{prop:reach} & --- & paper proof \\
Thm.~\ref{thm:decide} & --- & paper proof (reduction is Thm.~\ref{thm:char}) \\
Prop.~\ref{prop:coupling} & \texttt{coupling\_fails} & mechanised \\
Prop.~\ref{prop:degenerate} & \texttt{coupling}, \texttt{progComplete\_coupling} & mechanised \\
Prop.~\ref{prop:subsume} & \texttt{subsumes\_iff\_subset} & mechanised \\
Prop.~\ref{prop:factor} & \texttt{det\_factors} & mechanised (\texttt{rfl}) \\
Thm.~\ref{thm:suite} (finite) & \texttt{certify\_lb}, \texttt{certify\_ub} & mechanised (observed-set core) \\
Thm.~\ref{thm:suite} (infinite) & \texttt{certify\_infinite} & mechanised \\
Bridge & \texttt{realizable}, \texttt{progComplete\_iff\_complete} & mechanised \\
Prop.~\ref{prop:oceiling} (guard core) & \texttt{detW\_singleton\_of\_nonconstant} & mechanised \\
Prop.~\ref{prop:oconst} (identity) & \texttt{detW\_const} & mechanised (\texttt{rfl}) \\
Prop.~\ref{prop:pseudo} & \texttt{pseudo\_tested\_iff} & mechanised \\
Cor.~\ref{cor:relativise} (relativisation) & --- & paper proof \\
\bottomrule
\end{tabular}
\end{table}

Two caveats stated plainly. First, Theorem~\ref{thm:suite} is mechanised at the
level of the \emph{observed set}: \texttt{certify\_lb} and \texttt{certify\_ub}
establish that the minimum $|O|$ is $|I|$ for finite $I$, and the identity
$|T| = |O|$ for a minimal suite is the immediate preimage choice, done by hand.
Second, Theorem~\ref{thm:decide}'s spectrum is a meta-statement about
representations; its load-bearing reduction to footprint containment \emph{is}
the mechanised Theorem~\ref{thm:char}, but the complexity and computability
claims are paper proofs.

\paragraph{Artifact.}
The Lean development accompanies this paper. To reproduce: place the file in a
Lean~4 project with Mathlib, run \texttt{lake exe cache get}, and check e.g.\
\texttt{\#print axioms progComplete\_characterization}.

\section{Threats to Validity and Scope}
\label{sec:threats}

\paragraph{The exact oracle.}
Sections~\ref{sec:char}--\ref{sec:suite} assume the suite observes the full
return value. Section~\ref{sec:oracle} removes this assumption and shows the
theory relativises, with the ceiling unchanged for any non-constant oracle and
pseudo-testedness characterised (Proposition~\ref{prop:pseudo}). The relativised
results are paper proofs.

\paragraph{Purity and totality.}
Denotations are total functions. Void methods, exceptions, non-termination, and
side effects are outside the model. Programs with these features can be brought
inside by enlarging the return type with the corresponding observations, but we
do not develop the encoding, and the enlarged type is typically infinite, which
by Corollary~\ref{cor:infinite} puts absolute completeness out of reach and moves
the question to relative completeness.

\paragraph{Uniform post-composition.}
An output operator perturbs every returned value identically. A statement-level
mutation of one \texttt{return} site among several is not an output operator, and
the theory does not cover it.

\paragraph{Applicability of the ceiling.}
Theorem~\ref{thm:ceiling} shows a full absolute score is equivalent to output
coverage. This is a statement about what the score \emph{means}; it is not an
empirical claim about how often suites achieve it, nor about how well output
coverage correlates with fault detection. Those are separate, empirical
questions.

\paragraph{Scope of the negative results.}
Proposition~\ref{prop:coupling} shows coupling is not \emph{entailed}. It does
not show that coupling fails in practice, and it says nothing about program-text
mutation.

\section{Conclusion and Future Work}
\label{sec:conclusion}

Mutation testing is ordinarily used heuristically: a high score is read as
evidence of suite adequacy, justified by the coupling effect and the
competent-programmer hypothesis, which are empirical. For one operator class we
take a different stance---measuring a family not by the faults it can
\emph{express} but by what killing its mutants \emph{certifies}, uniformly over
programs and suites---and characterise the score exactly.

The characterisation partitions the content of a passing score into three
determinate parts: what it \textbf{certifies} (under an absolutely complete
family, precisely output coverage, Theorem~\ref{thm:ceiling}); what it
\textbf{provably cannot} certify (any fault distinguishing two inputs the program
maps to the same output, Remark~\ref{rem:ceiling}); and where the completeness
question is \textbf{undecidable} (relative completeness against a presented
target, off finite return types, Section~\ref{sec:relative}). For this class a
mutation score is an object with a characterised semantics rather than a quantity
trusted empirically.

Four directions remain open.

\begin{enumerate}
\item \textbf{Other operator classes.} Proposition~\ref{prop:factor} identifies
precisely what would let another class admit the same treatment: a
program-independent detection invariant analogous to the footprint.
Characterising the largest operator class whose detection factors through such an
invariant would say how far this style of result can reach.
\item \textbf{The weak-oracle theory.} Section~\ref{sec:oracle} shows the
development relativises and characterises pseudo-testedness. Mechanising the
relativised theory---abstracting the Lean development over the detection
function---and studying which oracles arise from real assertion styles is the
natural next step, and the one with the most contact with practice.
\item \textbf{Cost-optimal partial families.} Given a target class $\Gamma$ and a
per-operator cost, the minimum-cost $\Pi$ adequacy-complete for $\Gamma$ is a
weighted set-cover over footprints; Theorem~\ref{thm:char} makes the feasibility
predicate explicit.
\item \textbf{Decidable relative fragments.} Theorem~\ref{thm:decide} gives
polynomial time for explicit finite footprints and decidability for
Presburger-definable ones. A systematic account of the boundary---which footprint
representations admit decidable containment---is open.
\end{enumerate}

\begin{thebibliography}{99}

\bibitem{ammann2014}
Ammann, P., Delamaro, M.E., Offutt, J.: Establishing theoretical minimal sets of
mutants. In: ICST (2014)

\bibitem{budd1982}
Budd, T.A., Angluin, D.: Two notions of correctness and their relation to
testing. Acta Informatica \textbf{18}(1), 31--45 (1982)

\bibitem{demillo1978}
DeMillo, R.A., Lipton, R.J., Sayward, F.G.: Hints on test data selection: help
for the practicing programmer. IEEE Computer \textbf{11}(4), 34--41 (1978)

\bibitem{lean4}
de Moura, L., Ullrich, S.: The Lean 4 theorem prover and programming language.
In: CADE (2021)

\bibitem{goldman1995}
Goldman, S.A., Kearns, M.J.: On the complexity of teaching. J. Comput. Syst. Sci.
\textbf{50}(1), 20--31 (1995)

\bibitem{hamlet1977}
Hamlet, R.G.: Testing programs with the aid of a compiler. IEEE Trans. Softw.
Eng. \textbf{SE-3}(4), 279--290 (1977)

\bibitem{jia2011}
Jia, Y., Harman, M.: An analysis and survey of the development of mutation
testing. IEEE Trans. Softw. Eng. \textbf{37}(5), 649--678 (2011)

\bibitem{just2014}
Just, R., Jalali, D., Inozemtseva, L., Ernst, M.D., Holmes, R., Fraser, G.: Are
mutants a valid substitute for real faults in software testing? In: FSE (2014)

\bibitem{kurtz2015}
Kurtz, B., Ammann, P., Offutt, J.: Static analysis of mutant subsumption. In:
ICST Workshops (Mutation), pp. 1--10 (2015). \url{https://doi.org/10.1109/ICSTW.2015.7107454}

\bibitem{mathlib}
The mathlib Community: The Lean mathematical library. In: CPP (2020)

\bibitem{niedermayr2016}
Niedermayr, R., Juergens, E., Wagner, S.: Will my tests tell me if I break this
code? In: CSED@ICSE, pp. 23--29 (2016). \url{https://doi.org/10.1145/2896941.2896944}

\bibitem{offutt1992}
Offutt, A.J.: Investigations of the software testing coupling effect. ACM Trans.
Softw. Eng. Methodol. \textbf{1}(1), 5--20 (1992)

\bibitem{offutt1996}
Offutt, A.J., Lee, A., Rothermel, G., Untch, R.H., Zapf, C.: An experimental
determination of sufficient mutant operators. ACM Trans. Softw. Eng. Methodol.
\textbf{5}(2), 99--118 (1996)

\bibitem{papadakis2019}
Papadakis, M., Kintis, M., Zhang, J., Jia, Y., Le~Traon, Y., Harman, M.:
Mutation testing advances: an analysis and survey. Advances in Computers
\textbf{112}, 275--378 (2019)

\bibitem{veraperez2018}
Vera-P\'erez, O.L., Danglot, B., Monperrus, M., Baudry, B.: A comprehensive
study of pseudo-tested methods. Empirical Software Engineering \textbf{23}(3),
1908--1935 (2018)

\end{thebibliography}

\end{document}
